\usepackage{nomencl}
\makenomenclature

\providetoggle{nomsort}
\settoggle{nomsort}{true} % true = sort by use, false = sort as usual

\makeatletter
\iftoggle{nomsort}{%
    \let\old@@@nomenclature=\@@@nomenclature        
        \newcounter{@nomcount} \setcounter{@nomcount}{0}%
        \renewcommand\the@nomcount{\two@digits{\value{@nomcount}}}% Ensure 10>01
        \def\@@@nomenclature[#1]#2#3{% Taken from package documentation
          \addtocounter{@nomcount}{1}%
        \def\@tempa{#2}\def\@tempb{#3}%
          \protected@write\@nomenclaturefile{}%
          {\string\nomenclatureentry{\the@nomcount\nom@verb\@tempa @[{\nom@verb\@tempa}]%
          \begingroup\nom@verb\@tempb\protect\nomeqref{\theequation}%
          |nompageref}{\thepage}}%
          \endgroup
          \@esphack}%
      }{}
\makeatother

\nomenclature{\(\ZZ\)}{The integers.}
\nomenclature{\(\NN\)}{The natural numbers, usually starting at zero.}
\nomenclature{\(\ZZ_{\ge 0}\)}{The natural numbers, definitely starting at zero.}
\nomenclature{\(\ZZ_{\ge 1}\)}{The natural numbers, definitely starting at one.}
\nomenclature{\(\Meas{X}\)}{The set of measures of a measurable space \(X\), with the weak-\(\ast\) topology.}
\nomenclature{\(\Meas[1]{X}\)}{The set of probability measures of a measurable space \(X\), with the weak-\(\ast\) topology.}
\nomenclature{\(\mu_n \weakto \mu\)}{Weak-\(\ast\) convergence.}
\nomenclature{\(L^p(\mu)\)}{The space of \(p\)-norm integrable functions for \(\mu\) modulo the \almostevery equality relation.}
\nomenclature{\(L^p(\mathcal{A})\)}{The subspace of \(L^p(\mu)\) generated by functions that are \(\mathcal{A}\)-measurable.}
\nomenclature{\(\CE[\mu]{\phi, \mathcal{A}}\)}{Conditional expectation of \(\phi\) with respect to the measure \(\mu\) and \(\sigma\)-algebra \(\mathcal{A}\).}
\nomenclature{\((\Omega, \mathcal{F}, \Prob)\)}{Some arbitrary probability space.}
